
\documentclass[10pt]{article}

% Paquetes basicos
\usepackage[utf8]{inputenc}
\usepackage[spanish]{babel}
\usepackage{graphicx}
\usepackage{amsmath}
\usepackage{geometry}

% Margenes
\geometry{margin=1.5cm}

\begin{document}

% ENCABEZADO
\begin{center}
{\Large \textbf{REPORTE TECNICO}}\\
\textbf{Reto de Avionica AAIAQ 2026}\\
Ruta 1: Telemetria y adquisicion de datos
\end{center}

\noindent
\textbf{Nombre:} MOISES RICO\\
\textbf{Carrera:} TECNOLOGIA\\
\textbf{Semestre:} 1ERO

% OBJETIVO
\section{Objetivo}

Desarrollar un programa en Python para simular
el vuelo de un cohete y procesar los datos de
telemetria, identificando su altitud, velocidad
y los principales eventos durante el vuelo.

% DESARROLLO
\section{Desarrollo}

\subsection{Simulacion del vuelo}

Se programaron las siguientes fases:

\begin{itemize}
\item Rampa: el cohete permanece en reposo.
\item Propulsion: comienza el ascenso.
\item Vuelo balistico: actua la gravedad.
\item Paracaidas: comienza el descenso.
\item Aterrizaje: el cohete llega al suelo.
\end{itemize}

Para calcular la altitud se utilizo la siguiente
ecuacion de movimiento:

\begin{equation}
h_f=h_i+v_i\Delta t+
\frac{1}{2}a(\Delta t)^2
\end{equation}

\subsection{Procesamiento de datos}

Se utilizo Python para leer los datos del archivo
CSV y guardarlos en listas.

Despues se aplico un promedio movil para
suavizar las mediciones:

\begin{equation}
\bar{h}_i=\frac{1}{N}
\sum_{k=0}^{N-1}h_{i-k}
\end{equation}

Tambien se calculo la velocidad vertical:

\begin{equation}
v_i=\frac{h_i-h_{i-1}}{t_i-t_{i-1}}
\end{equation}

Finalmente, se programo la deteccion del
despegue, apogeo y aterrizaje.

% RESULTADOS
\section{Resultados}

Se generaron graficas para observar el
comportamiento del cohete durante el vuelo.

\begin{center}
\includegraphics[width=0.85\linewidth]
{comparacion_graficas.png}

Figura 1. Comparacion de altitud original y filtrada.
\end{center}

\begin{center}
\includegraphics[width=0.85\linewidth]
{velocidad_cohete.png}

Figura 2. Velocidad vertical del cohete.
\end{center}

\begin{center}
\includegraphics[width=0.85\linewidth]
{eventos_vuelo.png}

Figura 3. Eventos principales del vuelo.
\end{center}

Como se observa en las graficas, el filtrado
permite suavizar las mediciones y facilita
el analisis de la trayectoria.

% CONCLUSIONES
\section{Conclusiones}

El proyecto permitio aprender a utilizar Python
para simular un vuelo y procesar datos.

Se practicaron ciclos, condiciones, listas,
archivos CSV y graficas.

Como mejora futura se propone agregar ruido
a las mediciones y mejorar la limpieza de los
datos recibidos.

% BITACORA
\section{Bitacora de obstaculos}

Durante el desarrollo se presentaron algunas
dificultades:

\begin{itemize}

\item Se presentaron dificultades al crear y
organizar carpetas en GitHub. Se resolvieron
siguiendo los pasos de creacion de archivos.

\item Al comparar las graficas aparecieron
errores en el codigo. Se revisaron las
instrucciones y las variables utilizadas.

\item Se tuvieron dificultades para organizar
el reporte en LaTeX. Se utilizo Overleaf para
aprender a escribir y compilar el documento.

\end{itemize}

% REFERENCIAS
\section{Referencias}

\begin{itemize}
\item Documentacion de Python:
\texttt{https://docs.python.org/3/}

\item Documentacion de Matplotlib:
\texttt{https://matplotlib.org/}

\item Reto de Reclutamiento de Avionica AAIAQ 2026.

\item OpenAI ChatGPT: apoyo para comprender,
revisar y organizar el codigo y el reporte.
\end{itemize}

\end{document}
